\section{Knowledge gained}
\subsection{How EVA works now}
EVA is a local voice application with one visible character in Unity. After
loading, the user chooses a design, body-animation style, language and voice.
The six designs share the same conversation system. Extra scene hosts, if
present, are disabled during Play; three visible agents are not required.

The \textbf{backend} is the Python program behind Unity. A \textbf{WebSocket}
is the local connection carrying recordings, text, audio and mouth timings.
The \textbf{GPU} is the graphics card; it runs the models and draws the scene.

\begin{figure}[ht]
\centering
\begingroup\footnotesize
\renewcommand{\arraystretch}{1.25}
\begin{tabularx}{0.96\linewidth}{@{}P{7mm}Y@{}}
\toprule
1 & \textbf{Unity records speech} while B or the VR recording button is held.\\
2 & \textbf{Python receives the recording} with the selected language and voice profile.\\
3 & \textbf{Faster-Whisper creates text} using the selected language explicitly.\\
4 & \textbf{Ollama runs Llama 3.2} with a language instruction for this answer.\\
5 & \textbf{Piper creates speech and timings} with the selected matching voice.\\
6 & \textbf{Unity plays audio and animates EVA}, including mouth shapes and body states.\\
\bottomrule
\end{tabularx}
\endgroup
\caption{The current path from a question to a spoken answer.}
\label{fig:pipeline}
\end{figure}

Unity sends small audio pieces while the recording button is held. Whisper
returns provisional text, which may correct recent words. Releasing the button
finishes recognition and starts the answer. Answer text and audio arrive in
small pieces, so playback can begin before the whole answer is ready. Both text
boxes reveal complete words. Desktop text starts hidden and can be shown with T;
VR text starts visible and can be shown or hidden with left X.

\subsection{Tools and reasons for using them}
{\footnotesize
\begin{tabularx}{\linewidth}{@{}P{33mm}Y@{}}
\toprule
Tool & Job and reason\\
\midrule
Unity and C\# & Character display, microphone capture, menu, desktop/VR controls,
animation, audio and text boxes in one application.\\
Python & Connects the speech tools, manages conversation steps and sends results to Unity.\\
Faster-Whisper / CTranslate2 & Local speech recognition. Whisper Small runs on
CUDA with FP16 to use the NVIDIA card.\\
Ollama / Llama 3.2 3B & Ollama loads the local language model; Llama writes
answers. The 3B model limits memory use compared with larger models.\\
Piper / ONNX Runtime & Local English/German speech. ONNX Runtime runs the voice
models on CUDA; Piper's phoneme timings support lip-sync.\\
Microsoft Rocketbox & Existing rigged avatars, body animations and facial shapes,
including the 15 visemes used by EVA.\\
Input System / NativeWebSocket & Input handling and the persistent Unity--Python connection.\\
OpenXR / Quest Link & PC VR with Meta Quest 3 while Unity and models run on Windows.\\
\bottomrule
\end{tabularx}}
These choices support the current pipeline. Official tool links are in Appendix B.

\subsection{Menu, language and voice}
The menu provides six character designs and three body styles: Neutral,
Friendly and Expressive. The styles change body animation, not the language
model's personality. The states remain Idle, Listening, Thinking and Speaking.

\textbf{03 LANGUAGE fixes the whole turn.} English/Deutsch is sent to Whisper,
the answer instruction and the voice selector. There is no automatic language
detection or switching. Speak in the selected language and reopen settings
with M, or right B in VR, to change it while idle.

\textbf{04 VOICE is independent of the design.} English offers Kristin/female
and Mike/male; German offers Ramona/female and Thorsten/male. Clicking a name
selects it. Its small play icon only previews that voice. Loading and stop
symbols replace the play icon while needed. Preview audio uses a separate
AudioSource, does not alter the conversation and does not animate the mouth.
It stops on language/voice changes, conversation start or menu closure.

\subsection{VR interface and controls}
In VR, screen panels close to the eyes were hard to read at the edges of the
view. The interface now uses \textbf{World Space}: the menu sits beside EVA,
conversation text sits below the face, and loading appears in front of the user.
The surfaces are usually about 2--3 m away. They keep their position in the scene
when the user turns their head, so they can be viewed by looking towards EVA.

Use the right stick and A to select menu controls. Right B opens settings while
idle or closes the menu with Start/Continue. During conversation, hold right A
to record and release it to send. Release A after confirming a menu choice first.
The left stick moves and the right stick snap-turns by 30 degrees. Left X toggles
conversation text; left Y repositions the surfaces from the current position.
Opening the menu also repositions it. Head tracking remains active during menus,
while recording and walking stay blocked until loading and selection finish.

Distance, width and offsets can be changed in
\file{Assets/_Project/Resources/EVAControlSettings.asset}, under
\menu{VR panels (metres)}. Press Y after changing values during Play. Save the
asset outside Play to keep changes. Desktop layout and VR text visibility are
stored separately. Check readability in both eyes and while standing or seated.

\subsection{Problems encountered and lessons learned}
{\footnotesize
\begin{longtable}{@{}P{27mm}P{51mm}P{61mm}@{}}
\toprule
Area & What went wrong & Current solution or lesson\\
\midrule
\endhead
GPU dependencies & A CUDA DLL was missing; a CPU-only ONNX package also replaced
the GPU runtime. & Pin the NVIDIA packages, protect the GPU runtime in setup and
use the shared DLL configuration in \file{cuda_runtime.py}. Run preflight.\\
PC transfer & A copied environment pointed to the previous PC's Python and
exited with code 103. & Install Python 3.11 x64 locally and recreate only
\file{.venv}. Existing models and voices can stay.\\
Language choice & Automatic recognition, answer language and voice could
disagree. & Use one explicit menu language for all three stages. Piper itself
does not detect a text's language.\\
Complete playback & Speech ended early or filled the queue. & Wait for audio
playback to finish. The remaining 60-second queue limit is documented separately.\\
Lip-sync & The wrong face mesh had no visemes; mouth motion was too strong.
& Update face references when changing design and keep the scene's maximum
viseme weight at 50\%.\\
Voice vocabulary & Ramona's older model lacks some symbols produced by newer
eSpeak. & \file{piper_compat.py} filters unsupported phonemes before model input
and alignment. A regression test checks supported input is preserved.\\
Rig and turning & Bone axes differ; full-body turns could move EVA away.
& Keep import-time bone fixes, calibrated head axes and leg-only turn animation.\\
Server lifecycle & An old backend could remain after a script reload and fail
new connections. & Clean up owned process trees on disable/reload; wait for
\file{server.ready} and retain server tracebacks.\\
Preview UI & A new play icon lacked a CanvasRenderer and failed during raycasts.
& Require the component and add it explicitly before creating the icon.
Both creation paths and all icon states passed an isolated Unity test.\\
\bottomrule
\end{longtable}}

\subsection{Lip-sync and warm-up}
A \textbf{phoneme} is a speech sound; a \textbf{viseme} is a visible mouth shape.
Rocketbox provides 14 speech groups plus silence. Piper supplies phoneme timings,
Python maps them to these shapes, and Unity blends them against the audio clock.
No new lip-sync model was trained.

Warm-up prepares the first model calls: Whisper processes silence, Ollama writes
a short reply and the four Piper voices synthesize sample text. This loads model
resources and checks GPU execution; it does not train or fine-tune the models.
Input stays blocked until startup and the selection step are complete.

